\chapter{Information Theory}\label{information-theory}
\chapterstatus{draft}

How much information does a message contain, how far can it be compressed, and how fast can it be sent through a noisy channel? Shannon answered all three questions in one paper in 1948 \cite{shannon-1948}, with one
quantity, the \emph{entropy}, and the answers are sharp theorems: a source of entropy $H$ can be compressed to $H$ bits per symbol and no further, and a channel of capacity $C$ carries $C$ bits per use with arbitrarily small
error, and no more.
\begin{itemize}
\item Section~\ref{information-theory:section-entropy}: entropy, conditional entropy, mutual information; the data processing and Fano inequalities.
\item Section~\ref{information-theory:section-source-coding}: prefix codes, the Kraft--McMillan inequality, the source coding theorem and Huffman codes.
\item Section~\ref{information-theory:section-aep}: the asymptotic equipartition property and typical sequences.
\item Section~\ref{information-theory:section-channel-coding}: channels, capacity and Shannon's noisy channel coding theorem.
\item Section~\ref{information-theory:section-error-correcting}: error-correcting codes: Hamming, Reed--Solomon, and the bounds of Hamming, Singleton and Gilbert--Varshamov.
\item Section~\ref{information-theory:section-differential-entropy}: differential entropy and the Gaussian channel.
\item Section~\ref{information-theory:section-quantum}: a survey of quantum information, with the no-cloning theorem.
\end{itemize}

\noindent\textbf{Prerequisites.} Chapter~\ref{probability} (Probability Theory), Chapter~\ref{statistics} (Mathematical Statistics), for the Kullback--Leibler divergence, Chapter~\ref{fields} (Fields) and
Chapter~\ref{linear-algebra} (Linear Algebra).

\section{Entropy and mutual information}\label{information-theory:section-entropy}

\emph{Motivation.} How many yes/no questions are needed to identify an unknown object? If it is one of $2^k$ equally likely objects, $k$ questions suffice (halve the set each time) and are necessary. If the objects have
unequal probabilities, one should ask questions that split the probability, not the number of objects, in halves; an object of probability $p$ is then identified after about $\log_2\frac1p$ questions. The average
$\sum_xp(x)\log_2\frac{1}{p(x)}$ is the entropy.

\begin{definition}\label{information-theory:definition-entropy}
Let $X$ be a random variable with values in a finite or countable set $\mathcal{X}$ and probability mass function $p(x) = P(X = x)$. Its \emph{entropy} is
\[
H(X) = -\sum_{x \in \mathcal{X}}p(x)\log p(x) = E\Bigl[\log\frac{1}{p(X)}\Bigr] \in [0, \infty],
\]
with the convention $0\log 0 = 0$. For $p \in [0, 1]$, the \emph{binary entropy function} is $h(p) = -p\log p - (1 - p)\log(1 - p)$, the entropy of a Bernoulli variable with parameter $p$.
\end{definition}

\begin{notation}\label{information-theory:notation-log}
In this chapter $\log = \log_2$, and entropy is measured in \emph{bits}; $\ln$ is the natural logarithm, and entropies computed with $\ln$ are in \emph{nats} ($1$ nat $= \log_2e$ bits). $H(X)$ depends only on the law of
$X$, and we also write $H(p)$ for a probability vector $p$. For several variables, $p(x, y) = P(X = x, Y = y)$, $p(y \mid x) = P(Y = y \mid X = x)$, and $H(X, Y)$ is the entropy of the pair. $|\mathcal{X}|$ is the number of elements of
$\mathcal{X}$. $K(q, r)$ is the Kullback--Leibler divergence of Definition~\ref{statistics:definition-kullback-leibler}, always computed with the natural logarithm, so that $K/\ln 2$ is in bits.
\end{notation}

\begin{example}\label{information-theory:example-entropy}
\begin{enumerate}
\item A fair coin has $H = 1$ bit; a fair die $\log 6 \approx 2.585$ bits; a uniform variable on $N$ values $\log N$ bits.
\item $h(p)$ is symmetric about $\frac12$, concave, $h(0) = h(1) = 0$ and $h(\frac12) = 1$; $h(0.11) \approx 0.5$: a coin with $P(\text{heads}) = 0.11$ carries half a bit per toss.
\item If $X$ is geometric with $P(X = k) = 2^{-k}$, $k \geq 1$, then $H(X) = \sum_kk2^{-k} = 2$ bits. Questions ``is $X = 1$?'', ``is $X = 2$?'', \ldots{} identify $X$ with exactly $2$ questions on average.
\item Entropy can be infinite: $p(k) = \frac{c}{k\log^2k}$, $k \geq 2$, has $\sum p(k)\log\frac{1}{p(k)} = \infty$, by comparison with $\sum\frac{1}{k\log k}$.
\end{enumerate}
\end{example}

\begin{proposition}\label{information-theory:proposition-entropy-bounds}
$0 \leq H(X) \leq \log|\mathcal{X}|$ for finite $\mathcal{X}$. $H(X) = 0$ iff $X$ is almost surely constant, and $H(X) = \log|\mathcal{X}|$ iff $X$ is uniform.
\end{proposition}

\begin{proof}
Each term $-p\log p$ is $\geq 0$, and vanishes iff $p \in \{0, 1\}$. For the upper bound let $u$ be the uniform law; then
$K(p, u) = \sum_xp(x)\ln\frac{p(x)}{1/|\mathcal{X}|} = (\log|\mathcal{X}| - H(p))\ln 2$, which is $\geq 0$, with equality iff $p = u$, by Gibbs' inequality (Lemma~\ref{statistics:lemma-gibbs}).
\end{proof}

Under constraints on the mean the maximiser is not uniform but a Boltzmann distribution (Example~\ref{real-analysis:example-gibbs}).

\begin{definition}\label{information-theory:definition-conditional-entropy}
For random variables $X$, $Y$ with countable values, the \emph{conditional entropy} is
\[
H(Y \mid X) = \sum_xp(x)H(Y \mid X = x) = -\sum_{x,y}p(x, y)\log p(y \mid x),
\]
where $H(Y \mid X = x)$ is the entropy of the conditional law of $Y$ given $X = x$. The \emph{mutual information} is $I(X; Y) = H(X) - H(X \mid Y)$. For a third variable $Z$, $I(X; Y \mid Z) = H(X \mid Z) - H(X \mid Y, Z)$.
\end{definition}

\begin{proposition}\label{information-theory:proposition-entropy-rules}
For random variables with finite values:
\begin{enumerate}
\item (Chain rule) $H(X, Y) = H(X) + H(Y \mid X)$; more generally $H(X_1, \ldots, X_n) = \sum_iH(X_i \mid X_1, \ldots, X_{i-1})$.
\item $I(X; Y) = H(X) + H(Y) - H(X, Y) = I(Y; X) = K(p_{XY}, p_X \otimes p_Y)/\ln 2 \geq 0$, with equality iff $X$ and $Y$ are independent.
\item (Conditioning reduces entropy) $H(X \mid Y) \leq H(X)$, and $H(X_1, \ldots, X_n) \leq \sum_iH(X_i)$, with equality iff the $X_i$ are independent.
\item $I(X; Y \mid Z) \geq 0$, and (chain rule) $I(X; (Y, Z)) = I(X; Z) + I(X; Y \mid Z)$.
\item $H(f(X)) \leq H(X)$ for every function $f$, and $H(X \mid Y) = 0$ iff $X$ is almost surely a function of $Y$.
\end{enumerate}
\end{proposition}

\begin{proof}
(1) $\log p(x, y) = \log p(x) + \log p(y \mid x)$ where $p(x, y) > 0$; take expectations. Induction gives the general form.
(2) By (1) applied in both orders, $H(X) - H(X \mid Y) = H(X) + H(Y) - H(X, Y)$, which is symmetric, and equals $\sum_{x,y}p(x, y)\log\frac{p(x, y)}{p(x)p(y)}$. Gibbs' inequality gives the sign and the case of equality.
(3) is (2), and by induction with (1), $H(X_1, \ldots, X_n) = \sum_iH(X_i \mid X_1, \ldots, X_{i-1}) \leq \sum_iH(X_i)$; equality forces each $X_i$ to be independent of $(X_1, \ldots, X_{i-1})$.
(4) $I(X; Y \mid Z) = \sum_zp(z)I(X; Y \mid Z = z)$, a mixture of mutual informations of conditional laws, each $\geq 0$ by (2). The chain rule: $H(X) - H(X \mid Y, Z) = (H(X) - H(X \mid Z)) + (H(X \mid Z) - H(X \mid Y, Z))$.
(5) By (1), $H(X, f(X)) = H(X) + 0 = H(f(X)) + H(X \mid f(X)) \geq H(f(X))$. $H(X \mid Y) = 0$ iff each conditional law of $X$ given $Y = y$ with $p(y) > 0$ is a point mass (Proposition~\ref{information-theory:proposition-entropy-bounds}).
\end{proof}

\begin{counterexample}\label{information-theory:counterexample-conditioning}
Conditioning reduces entropy on average, not for each value: if $X \in \{0, 1\}$, $Y \in \{0, 1\}$ with $P(Y = 0) = \frac34$, $X = 0$ when $Y = 0$ and $X$ uniform when $Y = 1$, then $H(X) = h(\frac18) \approx 0.544$ but
$H(X \mid Y = 1) = 1$. Likewise $I(X; Y \mid Z)$ can exceed $I(X; Y)$: for independent fair bits $X$, $Y$ and $Z = X + Y \bmod 2$, $I(X; Y) = 0$ but $I(X; Y \mid Z) = 1$.
\end{counterexample}

\begin{theorem}[Data processing inequality]\label{information-theory:theorem-data-processing}
If $X \to Y \to Z$ is a \emph{Markov chain}, that is, $X$ and $Z$ are conditionally independent given $Y$, then $I(X; Z) \leq I(X; Y)$. In particular $I(X; g(Y)) \leq I(X; Y)$ for every function $g$.
\end{theorem}

\begin{proof}
By Proposition~\ref{information-theory:proposition-entropy-rules}(4) applied in two ways, $I(X; (Y, Z)) = I(X; Z) + I(X; Y \mid Z) = I(X; Y) + I(X; Z \mid Y)$, and $I(X; Z \mid Y) = 0$ by conditional independence (part (2)
applied to the conditional laws). Since $I(X; Y \mid Z) \geq 0$, $I(X; Z) \leq I(X; Y)$. For $Z = g(Y)$, $X \to Y \to g(Y)$ is a Markov chain.
\end{proof}

No processing of data, deterministic or random, can increase the information it contains about its source. In statistics, $T(Y)$ is sufficient for a parameter $\theta$ exactly when equality $I(\theta; T(Y)) = I(\theta; Y)$
holds for all priors (Section~\ref{statistics:section-models}, \cite{cover-thomas}).

\begin{theorem}[Fano's inequality]\label{information-theory:theorem-fano}
Let $X$ take values in a finite set $\mathcal{X}$, and let $\hat X = g(Y)$ be an estimate of $X$ from an observation $Y$, with error probability $P_e = P(\hat X \neq X)$. Then
\[
H(X \mid Y) \leq h(P_e) + P_e\log(|\mathcal{X}| - 1) \leq 1 + P_e\log|\mathcal{X}|.
\]
\end{theorem}

\begin{proof}
Let $E = \mathbf{1}_{\{\hat X \neq X\}}$. By the chain rule in two orders, $H(E, X \mid \hat X) = H(X \mid \hat X) + H(E \mid X, \hat X) = H(X \mid \hat X)$, as $E$ is a function of $(X, \hat X)$; and
$H(E, X \mid \hat X) = H(E \mid \hat X) + H(X \mid E, \hat X) \leq h(P_e) + P_e\log(|\mathcal{X}| - 1)$, because $H(E \mid \hat X) \leq H(E) = h(P_e)$, given $E = 0$ the variable $X = \hat X$ is determined, and given $E = 1$ it takes at most
$|\mathcal{X}| - 1$ values. Finally $H(X \mid Y) \leq H(X \mid \hat X)$, by the data processing inequality for $X \to Y \to \hat X$.
\end{proof}

\begin{remark}\label{information-theory:remark-axioms}
Entropy is not an arbitrary choice: up to the choice of unit, it is the only function of probability vectors that is continuous, symmetric, maximal at the uniform law, and satisfies the grouping rule
$H(p_1, \ldots, p_n) = H(p_1 + p_2, p_3, \ldots, p_n) + (p_1 + p_2)H(\frac{p_1}{p_1 + p_2}, \frac{p_2}{p_1 + p_2})$ (Shannon, Khinchin \cite{cover-thomas}). In statistical mechanics, Boltzmann's $S = k\ln W$ is the entropy of the uniform law
on $W$ microstates, in nats times Boltzmann's constant; in ergodic theory, the Kolmogorov--Sinai entropy of a dynamical system is the entropy per unit time of its best finite description (Chapter~\ref{ergodic-theory}).
\end{remark}

\section{Source coding}\label{information-theory:section-source-coding}

\begin{definition}\label{information-theory:definition-code}
Let $\mathcal{D}$ be a finite \emph{code alphabet} with $D \geq 2$ letters and $\mathcal{D}^*$ the set of finite words. A \emph{code} for a countable set $\mathcal{X}$ is a map $C \colon \mathcal{X} \to \mathcal{D}^*$, with \emph{codeword
lengths} $\ell(x) = |C(x)|$; its \emph{extension} sends a word $x_1\cdots x_n$ to the concatenation $C(x_1)\cdots C(x_n)$. The code is \emph{uniquely decodable} if its extension is injective, and a \emph{prefix code} if no
codeword is a prefix of another. For a random $X$, the \emph{expected length} is $L(C) = E[\ell(X)]$.
\end{definition}

A prefix code is uniquely decodable, and decodable \emph{instantaneously}: reading the encoded string from the left, a codeword is recognised as soon as it ends. Prefix codes over $\mathcal{D}$ are the same as sets of leaves of
the $D$-ary tree.

\begin{example}\label{information-theory:example-codes}
For $\mathcal{X} = \{a, b, c, d\}$ with probabilities $\frac12, \frac14, \frac18, \frac18$, the prefix code $a \mapsto 0$, $b \mapsto 10$, $c \mapsto 110$, $d \mapsto 111$ has expected length $1.75 = H(X)$. The code $a \mapsto 0$,
$b \mapsto 01$, $c \mapsto 011$, $d \mapsto 0111$ is uniquely decodable (every codeword starts with $0$, which marks the boundaries) but not prefix. The code $a \mapsto 0$, $b \mapsto 1$, $c \mapsto 01$ is not uniquely decodable:
$01$ is $c$ or $ab$. Morse code needs pauses between letters for the same reason.
\end{example}

\begin{theorem}[Kraft--McMillan inequality]\label{information-theory:theorem-kraft}
\begin{enumerate}
\item The lengths of a uniquely decodable code satisfy $\sum_xD^{-\ell(x)} \leq 1$.
\item Conversely, if positive integers $\ell(x)$ satisfy $\sum_xD^{-\ell(x)} \leq 1$, there is a prefix code with these lengths.
\end{enumerate}
\end{theorem}

\begin{proof}
(1) \emph{Finite case.} Let $\mathcal{X}$ be finite, $\ell_{\max} = \max\ell(x)$ and $S = \sum_xD^{-\ell(x)}$. Then $S^k = \sum_{x_1, \ldots, x_k}D^{-(\ell(x_1) + \cdots + \ell(x_k))} = \sum_{m=1}^{k\ell_{\max}}a_mD^{-m}$, where $a_m$ is the number of words of
length $k$ in $\mathcal{X}$ whose encoding has length $m$. By unique decodability these encodings are distinct words of length $m$, so $a_m \leq D^m$, and $S^k \leq k\ell_{\max}$. As $k \to \infty$, $S \leq (k\ell_{\max})^{1/k} \to 1$. \emph{Countable
case.} Every finite subcode is uniquely decodable, so every partial sum is at most $1$.
(2) Enumerate $\mathcal{X}$ as $x_1, x_2, \ldots$ with $\ell_1 \leq \ell_2 \leq \cdots$ (possible since each value $\ell$ is taken by at most $D^\ell$ elements, by the inequality). Let $s_i = \sum_{j < i}D^{-\ell_j} \leq 1 - D^{-\ell_i} < 1$. As $\ell_j \leq \ell_i$ for $j < i$, $s_i$ is a multiple of $D^{-\ell_i}$, so its base-$D$ expansion terminates within $\ell_i$ digits; let $C(x_i)$ be its
first $\ell_i$ digits (padded with zeros). For $j > i$, $s_j \geq s_i + D^{-\ell_i}$, so $s_j \notin [s_i, s_i + D^{-\ell_i})$, the set of numbers whose expansion starts with $C(x_i)$; hence
$C(x_i)$ is not a prefix of $C(x_j)$, and it is not a prefix of an earlier, not longer, codeword either, being different from it.
\end{proof}

So nothing is lost by restricting to prefix codes: any lengths achievable by a uniquely decodable code are achievable by a prefix code.

\begin{theorem}[Source coding theorem]\label{information-theory:theorem-source-coding}
Let $H_D(X) = H(X)/\log D$ be the entropy in base $D$.
\begin{enumerate}
\item Every uniquely decodable code has $L(C) \geq H_D(X)$, with equality iff $p(x) = D^{-\ell(x)}$ for all $x$.
\item There is a prefix code with $L(C) < H_D(X) + 1$, for instance with the \emph{Shannon lengths} $\ell(x) = \lceil\log_D\frac{1}{p(x)}\rceil$.
\item Coding blocks of $n$ i.i.d.\ copies of $X$, the optimal expected length per symbol $L_n^*/n$ satisfies $H_D(X) \leq L_n^*/n < H_D(X) + \frac1n$, and tends to $H_D(X)$.
\end{enumerate}
\end{theorem}

\begin{proof}
(1) Let $r(x) = D^{-\ell(x)}/S$ with $S = \sum D^{-\ell(x)} \leq 1$ (Theorem~\ref{information-theory:theorem-kraft}). Then
$L(C) - H_D(X) = \sum_xp(x)\log_D\frac{p(x)}{D^{-\ell(x)}} = \frac{K(p, r)}{\ln D} - \log_DS \geq 0$, by Gibbs' inequality and $S \leq 1$, with equality iff $p = r$ and $S = 1$.
(2) The Shannon lengths satisfy $D^{-\ell(x)} \leq p(x)$, so Kraft's inequality holds and a prefix code exists; and $\ell(x) < \log_D\frac{1}{p(x)} + 1$.
(3) Apply (1) and (2) to $(X_1, \ldots, X_n)$, whose entropy is $nH(X)$ (Proposition~\ref{information-theory:proposition-entropy-rules}(3)).
\end{proof}

\begin{theorem}[Huffman]\label{information-theory:theorem-huffman}
For a finite $\mathcal{X}$ with probabilities $p_1 \geq p_2 \geq \cdots \geq p_m > 0$ ($m \geq 2$), the binary prefix code built recursively by merging the two least probable symbols into one symbol of probability $p_{m-1} + p_m$,
coding the reduced source, and appending $0$ and $1$ to the codeword of the merged symbol, has minimal expected length among all uniquely decodable binary codes.
\end{theorem}

\begin{proof}
By Theorem~\ref{information-theory:theorem-kraft} it suffices to compare with prefix codes, and an optimal prefix code exists (finitely many length vectors have expected length below a given bound). \emph{An optimal code with
good shape.} In an optimal prefix code: if $p_i > p_j$ then $\ell_i \leq \ell_j$ (otherwise swapping the codewords shortens the code); the longest codewords come in sibling pairs (a longest codeword without a sibling could be
shortened by deleting its last digit, keeping the prefix property); so after swapping codewords of equal length we may assume that $x_{m-1}$ and $x_m$ are siblings of maximal length. \emph{Induction.} For such a code $C$,
merging $x_{m-1}$ and $x_m$ gives a prefix code $C^{\prime}$ for the reduced source with $L(C) = L(C^{\prime}) + p_{m-1} + p_m$. Conversely the Huffman code $C_H$ of the source and the Huffman code $C_H^{\prime}$ of the reduced source satisfy
$L(C_H) = L(C_H^{\prime}) + p_{m-1} + p_m$. By induction on $m$, $L(C_H^{\prime}) \leq L(C^{\prime})$, so $L(C_H) \leq L(C)$.
\end{proof}

\begin{example}\label{information-theory:example-huffman}
For probabilities $0.4, 0.2, 0.2, 0.1, 0.1$, Huffman merges $0.1 + 0.1 = 0.2$, then $0.2 + 0.2 = 0.4$, then $0.2 + 0.4 = 0.6$, then $0.4 + 0.6$, giving lengths $1, 2, 3, 4, 4$ (one of several optimal choices) and expected length
$2.2$ bits, against $H \approx 2.122$. Huffman found the algorithm in 1951 as a student, in a term paper that replaced the final exam of Fano's course \cite{huffman-1952}.
\end{example}

\begin{remark}\label{information-theory:remark-practical-compression}
\emph{Arithmetic coding} encodes a whole message as a subinterval of $[0, 1)$ of length equal to its probability, and achieves $-\log p(x_1, \ldots, x_n) + 2$ bits for a message, with adaptive probability models. \emph{Universal}
methods such as the Lempel--Ziv algorithms (used in \texttt{zip} and \texttt{gzip}) need no model at all and still achieve the entropy rate of every stationary ergodic source asymptotically \cite{cover-thomas}. Lossy
compression (JPEG, MP3) is governed by Shannon's \emph{rate--distortion theory}: the minimal rate for reproducing a source within distortion $\Delta$ is $R(\Delta) = \min I(X; \hat X)$ over conditional laws with
$E[d(X, \hat X)] \leq \Delta$.
\end{remark}

\section{The asymptotic equipartition property}\label{information-theory:section-aep}

\begin{theorem}[Asymptotic equipartition property]\label{information-theory:theorem-aep}
Let $X_1, X_2, \ldots$ be i.i.d.\ with values in a finite set $\mathcal{X}$ and law $p$, and $p(x_1, \ldots, x_n) = \prod_ip(x_i)$. Then $-\frac1n\log p(X_1, \ldots, X_n) \to H(X_1)$ almost surely. For $\varepsilon > 0$ the \emph{typical set}
\[
A_\varepsilon^{(n)} = \bigl\{(x_1, \ldots, x_n) : 2^{-n(H + \varepsilon)} \leq p(x_1, \ldots, x_n) \leq 2^{-n(H - \varepsilon)}\bigr\}
\]
satisfies $P(A_\varepsilon^{(n)}) \to 1$, $|A_\varepsilon^{(n)}| \leq 2^{n(H + \varepsilon)}$ for all $n$, and $|A_\varepsilon^{(n)}| \geq (1 - \varepsilon)2^{n(H - \varepsilon)}$ for large $n$.
\end{theorem}

\begin{proof}
$-\frac1n\log p(X_1, \ldots, X_n) = \frac1n\sum_i\log\frac{1}{p(X_i)}$ is an average of i.i.d.\ bounded variables with mean $H$, so the strong law (Theorem~\ref{probability:theorem-strong-law}) gives the first claim, and
$P(A_\varepsilon^{(n)}) \to 1$ follows. Then $1 \geq \sum_{x \in A_\varepsilon^{(n)}}p(x) \geq |A_\varepsilon^{(n)}|2^{-n(H + \varepsilon)}$, and for large $n$, $1 - \varepsilon \leq P(A_\varepsilon^{(n)}) \leq |A_\varepsilon^{(n)}|2^{-n(H - \varepsilon)}$.
\end{proof}

So for large $n$ almost all the probability sits on about $2^{nH}$ sequences, each of probability about $2^{-nH}$: the sample behaves like a uniform choice among $2^{nH}$ possibilities. When $H < \log|\mathcal{X}|$ this is a
vanishing fraction of all $|\mathcal{X}|^n$ sequences. Note that the single most probable sequence is usually \emph{not} typical: for a coin with $P(1) = 0.9$, it is $11\cdots1$, while typical sequences contain about $10\%$ zeros.

\begin{corollary}[Fixed-length compression]\label{information-theory:corollary-fixed-length}
For every $R > H(X_1)$ and $\delta > 0$ there are, for large $n$, maps $f_n \colon \mathcal{X}^n \to \{0, 1\}^{\lceil nR\rceil}$ and $g_n$ in the other direction with $P(g_n(f_n(X_1, \ldots, X_n)) \neq (X_1, \ldots, X_n)) < \delta$. For
$R < H(X_1)$, every such pair has error probability tending to $1$.
\end{corollary}

\begin{proof}
For $R > H$ take $\varepsilon = R - H$ and encode the at most $2^{n(H + \varepsilon)}$ typical sequences injectively, everything else by a fixed word. For $R < H$, choose $\varepsilon$ with $R < H - 2\varepsilon$. The set $B$ of
correctly decoded sequences has at most $2^{\lceil nR\rceil}$ elements, so $P(B) \leq P(B \cap A_\varepsilon^{(n)}) + P(\text{not typical}) \leq 2^{nR + 1}2^{-n(H - \varepsilon)} + o(1) \to 0$.
\end{proof}

\section{Channel coding}\label{information-theory:section-channel-coding}

\begin{definition}\label{information-theory:definition-channel}
A \emph{discrete memoryless channel} consists of finite input and output alphabets $\mathcal{X}$, $\mathcal{Y}$ and transition probabilities $W(y \mid x)$; used $n$ times on inputs $x_1, \ldots, x_n$ it produces outputs with probability
$\prod_iW(y_i \mid x_i)$. Its \emph{capacity} is
\[
C = \max_{p_X}I(X; Y),
\]
the maximum over input laws $p_X$ of the mutual information, where $Y$ has conditional law $W(\cdot \mid x)$ given $X = x$ (the maximum exists: $I$ is continuous on the compact simplex of input laws). An $(M, n)$ \emph{code} consists of
an encoder $\{1, \ldots, M\} \to \mathcal{X}^n$ and a decoder $\mathcal{Y}^n \to \{1, \ldots, M\}$; its \emph{rate} is $R = \frac{\log M}{n}$ bits per channel use, and its \emph{average error probability} $P_e^{(n)}$ is the probability of
decoding wrongly when the message is uniform.
\end{definition}

\begin{example}\label{information-theory:example-capacity}
\begin{enumerate}
\item The \emph{binary symmetric channel} flips each bit with probability $p$. For any input law, $H(Y \mid X) = h(p)$, so $I(X; Y) = H(Y) - h(p) \leq 1 - h(p)$, with equality for the uniform input: $C = 1 - h(p)$. For $p = 0.11$,
$C \approx 0.5$; for $p = \frac12$, $C = 0$.
\item The \emph{binary erasure channel} delivers each bit with probability $1 - \alpha$ and replaces it by a symbol ``?'' otherwise. Then $C = 1 - \alpha$: the receiver learns which bits were lost, and the capacity is exactly the
fraction that arrives.
\item A \emph{noisy typewriter} on $26$ letters sends each letter to itself or to the next one with probability $\frac12$ each. Using only every other letter there is no error at all, at rate $\log 13$; and indeed
$C = \log 26 - 1 = \log 13$.
\end{enumerate}
\end{example}

\begin{theorem}[Shannon's noisy channel coding theorem]\label{information-theory:theorem-channel-coding}
Let $C$ be the capacity of a discrete memoryless channel.
\begin{enumerate}
\item (Achievability) For every $R < C$ there are $(\lceil 2^{nR}\rceil, n)$ codes with $P_e^{(n)} \to 0$ as $n \to \infty$.
\item (Converse) Every sequence of $(\lceil 2^{nR}\rceil, n)$ codes with $P_e^{(n)} \to 0$ has $R \leq C$; precisely, $P_e^{(n)} \geq 1 - \frac{C}{R} - \frac{1}{nR}$.
\end{enumerate}
\end{theorem}

\begin{proof}
(2) Let the message $V$ be uniform on $M = \lceil 2^{nR}\rceil$ values, $X^n$ its codeword, $Y^n$ the output and $\hat V$ the decoded message. By memorylessness, $H(Y^n \mid X^n) = \sum_iH(Y_i \mid X_i)$, so
$I(X^n; Y^n) = H(Y^n) - \sum_iH(Y_i \mid X_i) \leq \sum_i(H(Y_i) - H(Y_i \mid X_i)) = \sum_iI(X_i; Y_i) \leq nC$ (Proposition~\ref{information-theory:proposition-entropy-rules}(3)). Since $V \to X^n \to Y^n$ is a Markov chain, so is $Y^n \to X^n \to V$ (conditional independence is symmetric), and Fano's inequality for $\hat V$ together with the
data processing inequality $I(V; Y^n) \leq I(X^n; Y^n)$ give
\[\begin{aligned}
nR \leq \log M = H(V) &= H(V \mid Y^n) + I(V; Y^n)\\
&\leq 1 + P_e^{(n)}\log M + I(X^n; Y^n) \leq 1 + P_e^{(n)}\log M + nC.
\end{aligned}\]
Dividing by $\log M \geq nR$ gives the bound.
\end{proof}

\begin{proof}[Sketch of proof of (1)]
\emph{Random coding.} Fix an input law $p$ achieving $C$ and choose the $M$ codewords independently, each with i.i.d.\ coordinates of law $p$. \emph{Joint typicality decoding.} Decode $y^n$ as the unique message whose
codeword is jointly typical with $y^n$, that is, $(x^n, y^n)$ is typical for the joint law $p(x)W(y \mid x)$ in the sense of Theorem~\ref{information-theory:theorem-aep} (and its marginals are typical). \emph{Error analysis.} By the
AEP the transmitted codeword is jointly typical with the output with probability tending to $1$; a different, independent codeword is jointly typical with the output with probability at most $2^{-n(I(X; Y) - 3\varepsilon)}$, by
counting typical pairs. So the average error over the random choice of the code is at most $\varepsilon + M2^{-n(C - 3\varepsilon)}$, which tends to $0$ when $R < C - 3\varepsilon$, and some code does at least as well as the average.
The omitted verification is the joint typicality lemma \cite{cover-thomas}. Discarding the worse half of the codewords even gives small \emph{maximal} error probability.
\end{proof}

The theorem is an existence statement, and for half a century the codes it promises were out of reach. Turbo codes (1993), low-density parity-check codes and polar codes (Ar\i kan, 2009) now operate close to capacity, in
mobile phones and satellite links.

\begin{remark}\label{information-theory:remark-feedback}
Feedback does not increase the capacity of a memoryless channel (Shannon), although it simplifies coding; the capacity of a channel with memory, or of networks of channels, is in general unknown. For the binary symmetric
channel the error probability of the best codes decays exponentially in $n$ for every $R < C$, with an exponent known exactly at rates near $C$.
\end{remark}

\section{Error-correcting codes}\label{information-theory:section-error-correcting}

Shannon's theorem uses long random codes. The algebraic theory of codes constructs explicit codes with structure that makes encoding and decoding efficient.

\begin{definition}\label{information-theory:definition-block-code}
Let $\mathbb{F}_q$ be the finite field with $q$ elements (Theorem~\ref{fields:theorem-finite-fields}). The \emph{Hamming distance} on $\mathbb{F}_q^n$ is $d(x, y) = \#\{i : x_i \neq y_i\}$, and the \emph{weight} of $x$ is $\mathrm{wt}(x) = d(x, 0)$. A
\emph{code} of length $n$ is a subset $\mathcal{C} \subseteq \mathbb{F}_q^n$; its \emph{minimum distance} $d(\mathcal{C})$ is the least distance between distinct codewords. A \emph{linear code} $[n, k, d]_q$ is a $k$-dimensional subspace with
minimum distance $d$; a \emph{generator matrix} is a $k \times n$ matrix whose rows form a basis, and a \emph{parity check matrix} is an $(n - k) \times n$ matrix $H$ of rank $n - k$ with $\mathcal{C} = \ker H$.
\end{definition}

\begin{proposition}\label{information-theory:proposition-minimum-distance}
\begin{enumerate}
\item The Hamming distance is a metric. A code of minimum distance $d$ detects up to $d - 1$ errors and corrects up to $t = \lfloor\frac{d - 1}{2}\rfloor$ errors: the balls of radius $t$ around codewords are disjoint, so decoding to the
nearest codeword is correct when at most $t$ coordinates are changed.
\item For a linear code, $d(\mathcal{C}) = \min\{\mathrm{wt}(c) : c \in \mathcal{C}, c \neq 0\}$, and $d(\mathcal{C})$ is the least number of linearly dependent columns of a parity check matrix.
\end{enumerate}
\end{proposition}

\begin{proof}
(1) The triangle inequality holds coordinatewise. If $x$ lay in the balls of radius $t$ around codewords $c \neq c^{\prime}$, then $d(c, c^{\prime}) \leq 2t < d$. If at most $d - 1$ coordinates of a codeword are changed, the result is not
another codeword. (2) $d(c, c^{\prime}) = \mathrm{wt}(c - c^{\prime})$ and $c - c^{\prime} \in \mathcal{C}$. A nonzero $c$ with $Hc = 0$ of weight $w$ is a linear dependence among $w$ columns of $H$, and conversely.
\end{proof}

\begin{example}[Hamming codes]\label{information-theory:example-hamming-codes}
Let $r \geq 2$ and let $H$ be the $r \times (2^r - 1)$ binary matrix whose columns are all nonzero vectors of $\mathbb{F}_2^r$. Its kernel is the \emph{Hamming code} $[2^r - 1, 2^r - 1 - r, 3]_2$: no column is zero and no two
columns are equal, so no one or two columns are dependent, while $e_1$, $e_2$, $e_1 + e_2$ are. It corrects one error, and decoding is immediate: if one bit is flipped, the \emph{syndrome} $Hy$ is the column of $H$ at the
position of the error. For $r = 3$, the $[7, 4, 3]$ code sends $4$ data bits with $3$ check bits; Hamming invented it in 1950 because the relay computers at Bell Labs stopped at every error over weekends.
\end{example}

\begin{theorem}[Bounds on codes]\label{information-theory:theorem-code-bounds}
Let $\mathcal{C} \subseteq \mathbb{F}_q^n$ have $M$ codewords and minimum distance $d$, and $V_q(n, r) = \sum_{i=0}^r\binom ni(q - 1)^i$ the number of points in a Hamming ball of radius $r$.
\begin{enumerate}
\item (Hamming, sphere packing) $M\,V_q(n, t) \leq q^n$ with $t = \lfloor\frac{d - 1}{2}\rfloor$.
\item (Singleton) $M \leq q^{n - d + 1}$; for linear codes, $d \leq n - k + 1$.
\item (Gilbert--Varshamov) Conversely, for every $d \leq n$ there is a code of length $n$ and minimum distance at least $d$ with $M \geq q^n/V_q(n, d - 1)$.
\end{enumerate}
\end{theorem}

\begin{proof}
(1) The balls of radius $t$ around codewords are disjoint (Proposition~\ref{information-theory:proposition-minimum-distance}). (2) Deleting the last $d - 1$ coordinates is injective on $\mathcal{C}$, since distinct codewords differ in at
least $d$ places; its image lies in $\mathbb{F}_q^{n - d + 1}$. (3) Choose codewords greedily, each at distance $\geq d$ from the previous ones, until no point is left; then the balls of radius $d - 1$ around the $M$ codewords cover
$\mathbb{F}_q^n$, so $MV_q(n, d - 1) \geq q^n$.
\end{proof}

Codes attaining the Hamming bound are \emph{perfect}; the Hamming codes are perfect, as $2^{2^r - 1 - r}(1 + (2^r - 1)) = 2^{2^r - 1}$. Apart from trivial codes, the only other perfect codes over prime power alphabets are the binary
$[23, 12, 7]$ and ternary $[11, 6, 5]$ Golay codes (Tiet\"av\"ainen, van Lint) \cite{macwilliams-sloane}. Codes attaining the Singleton bound are \emph{maximum distance separable} (MDS).

\begin{theorem}[Reed--Solomon codes]\label{information-theory:theorem-reed-solomon}
Let $a_1, \ldots, a_n$ be distinct elements of $\mathbb{F}_q$ ($n \leq q$) and $1 \leq k \leq n$. The \emph{Reed--Solomon code} $\{(f(a_1), \ldots, f(a_n)) : f \in \mathbb{F}_q[x], \deg f < k\}$ is a linear $[n, k, n - k + 1]_q$ code; it is MDS.
\end{theorem}

\begin{proof}
The evaluation map $f \mapsto (f(a_i))_i$ is linear on the $k$-dimensional space of polynomials of degree $< k$. A nonzero such $f$ has at most $k - 1$ roots (Corollary~\ref{rings:corollary-roots}), so its image has weight at least
$n - k + 1$; in particular the map is injective, and the code has dimension $k$. By the Singleton bound, $d = n - k + 1$.
\end{proof}

\begin{example}\label{information-theory:example-reed-solomon}
Reed--Solomon codes over $\mathbb{F}_{256}$, whose symbols are bytes, correct bursts of errors: a scratch destroys many consecutive bits but only a few bytes. Compact discs use two interleaved Reed--Solomon codes of lengths
$32$ and $28$ over $\mathbb{F}_{256}$; QR codes, DVDs, deep-space probes and RAID storage use them too. They can be decoded efficiently up to $\lfloor\frac{n - k}{2}\rfloor$ errors by the Berlekamp--Welch or Berlekamp--Massey algorithms,
which reduce decoding to linear algebra over $\mathbb{F}_q$ \cite{macwilliams-sloane}.
\end{example}

\begin{remark}\label{information-theory:remark-asymptotic-codes}
For binary codes with relative distance $\delta = d/n$ and rate $R = k/n$, the Gilbert--Varshamov bound gives codes with $R \geq 1 - h(\delta)$, using $V_2(n, \delta n) \leq 2^{nh(\delta)}$ for $\delta \leq \frac12$. Whether better binary codes
exist is a famous open problem; for $q \geq 49$ a square, algebraic-geometry codes built from curves with many points over $\mathbb{F}_q$ beat the Gilbert--Varshamov bound (Tsfasman--Vl\u{a}du\c{t}--Zink), one of the striking
applications of algebraic geometry (Chapter~\ref{curves}).
\end{remark}

\section{Differential entropy and the Gaussian channel}\label{information-theory:section-differential-entropy}

\begin{definition}\label{information-theory:definition-differential-entropy}
For a real random variable with density $f$, the \emph{differential entropy} is $h(X) = -\int f\ln f\,dx$ (in nats), when the integral exists; for a random vector with density on $\RR^n$ the definition is the same.
\end{definition}

\begin{example}\label{information-theory:example-differential-entropy}
The uniform law on $[a, b]$ has $h = \ln(b - a)$, which is negative for $b - a < 1$: differential entropy is not a limit of discrete entropies, but $H$ of the quantisation of $X$ to intervals of length $\Delta$ is
$h(X) - \ln\Delta + o(1)$ for nice densities. $N(m, \sigma^2)$ has $h = \frac12\ln(2\pi e\sigma^2)$: $-\ln f = \frac12\ln(2\pi\sigma^2) + \frac{(x - m)^2}{2\sigma^2}$ has mean $\frac12\ln(2\pi\sigma^2) + \frac12$. Exponential: $h = 1 - \ln\lambda$.
\end{example}

\begin{proposition}[The Gaussian maximises entropy]\label{information-theory:proposition-gaussian-max-entropy}
If $X$ has density $f$ and variance $\sigma^2$, then $h(X) \leq \frac12\ln(2\pi e\sigma^2)$, with equality iff $X$ is normal.
\end{proposition}

\begin{proof}
Let $g$ be the $N(m, \sigma^2)$ density with $m = E[X]$. By Gibbs' inequality (Lemma~\ref{statistics:lemma-gibbs}, with respect to Lebesgue measure), $0 \leq K(f, g) = -h(X) - \int f\ln g$, and
$-\int f\ln g = \frac12\ln(2\pi\sigma^2) + \frac{E[(X - m)^2]}{2\sigma^2} = \frac12\ln(2\pi e\sigma^2)$. Equality holds iff $f = g$ almost everywhere.
\end{proof}

This explains the special role of Gaussian noise: for a given power, it is the most random and hence the worst noise. In the central limit theorem, the entropy of the normalised sums increases to that of the Gaussian
(Barron), a form of the second law of thermodynamics.

\begin{theorem}[Capacity of the Gaussian channel]\label{information-theory:theorem-gaussian-channel}
For the channel $Y = X + Z$ with Gaussian noise $Z \sim N(0, N)$ independent of the input, and input power constraint $\frac1n\sum_ix_i^2 \leq P$ on codewords, the capacity is
\[
C = \frac12\log\Bigl(1 + \frac PN\Bigr) \text{ bits per channel use}.
\]
\end{theorem}

\begin{proof}[Sketch of proof]
The coding theorem and its converse hold as in Theorem~\ref{information-theory:theorem-channel-coding}, with $C = \max\{I(X; Y) : E[X^2] \leq P\}$ and $I(X; Y) = h(Y) - h(Y \mid X) = h(Y) - h(Z)$. Since $\operatorname{Var}Y \leq P + N$,
Proposition~\ref{information-theory:proposition-gaussian-max-entropy} gives $h(Y) \leq \frac12\ln(2\pi e(P + N))$, with equality for $X \sim N(0, P)$; and $h(Z) = \frac12\ln(2\pi eN)$. The omitted verifications are the extension of the
coding theorem to continuous alphabets with a power constraint \cite{cover-thomas}.
\end{proof}

A band-limited channel of bandwidth $W$ hertz carries $2W$ independent samples per second (by the sampling theorem), which gives the \emph{Shannon--Hartley formula} $C = W\log(1 + \mathrm{SNR})$ bits per second. A telephone line
with $W = 3000$ Hz and signal-to-noise ratio $30$ dB ($\mathrm{SNR} = 1000$) has capacity about $30000$ bits per second, which is what the last generation of analogue modems achieved.

\begin{remark}\label{information-theory:remark-epi}
The \emph{entropy power inequality} of Shannon and Stam, $e^{2h(X + Y)/n} \geq e^{2h(X)/n} + e^{2h(Y)/n}$ for independent random vectors in $\RR^n$, with equality for proportional Gaussians, is a deep analogue of the
Brunn--Minkowski inequality of convex geometry; both follow from sharp forms of Young's convolution inequality \cite{cover-thomas}.
\end{remark}

\section{Quantum information}\label{information-theory:section-quantum}

In quantum mechanics the state of a system is a unit vector $\psi$ in a complex Hilbert space $\mathcal{H}$, up to a phase, or more generally a \emph{density operator} $\rho \geq 0$ with $\operatorname{tr}\rho = 1$ (Chapter~\ref{quantum-mechanics}).
A \emph{qubit} is a system with $\mathcal{H} = \CC^2$, with basis $|0\rangle$, $|1\rangle$; unlike a bit it can be in any superposition $\alpha|0\rangle + \beta|1\rangle$, $|\alpha|^2 + |\beta|^2 = 1$, and measuring it gives $0$ or $1$ with
probabilities $|\alpha|^2$, $|\beta|^2$. Two systems combine by the tensor product $\mathcal{H}_1 \otimes \mathcal{H}_2$, with the inner product $\langle a \otimes b, c \otimes d\rangle = \langle a, c\rangle\langle b, d\rangle$ on pure tensors, and
closed systems evolve by unitary operators.

\begin{theorem}[No-cloning theorem]\label{information-theory:theorem-no-cloning}
Let $\dim\mathcal{H} \geq 2$ and $e \in \mathcal{H}$ a unit vector. There is no unitary (or even isometric) operator $U$ on $\mathcal{H} \otimes \mathcal{H}$ with $U(\psi \otimes e) = \psi \otimes \psi$ for every unit vector $\psi \in \mathcal{H}$.
\end{theorem}

\begin{proof}
If $U$ preserves inner products, then for unit vectors $\phi$, $\psi$,
$\langle\phi, \psi\rangle = \langle\phi, \psi\rangle\langle e, e\rangle = \langle U(\phi \otimes e), U(\psi \otimes e)\rangle = \langle\phi \otimes \phi, \psi \otimes \psi\rangle = \langle\phi, \psi\rangle^2$, so $\langle\phi, \psi\rangle \in \{0, 1\}$. But for orthonormal $u$, $v$, the
unit vectors $\phi = u$ and $\psi = (u + v)/\sqrt2$ have $\langle\phi, \psi\rangle = \frac{1}{\sqrt2}$.
\end{proof}

Unknown quantum states cannot be copied, unlike classical bits; this is the basis of quantum cryptography (the BB84 protocol), where an eavesdropper cannot copy the transmitted qubits without disturbing them.

\begin{remark}\label{information-theory:remark-quantum-survey}
\begin{enumerate}
\item \emph{Von Neumann entropy.} $S(\rho) = -\operatorname{tr}\rho\log\rho = H(\lambda)$, the Shannon entropy of the eigenvalues of $\rho$. Schumacher's theorem is the quantum source coding theorem: a source of pure states with
average density operator $\rho$ can be compressed to $S(\rho)$ qubits per state.
\item \emph{Entanglement.} The \emph{Bell state} $\frac{1}{\sqrt2}(|00\rangle + |11\rangle) \in \CC^2 \otimes \CC^2$ is not a product of two vectors; its parts are each in the maximally mixed state $\frac12I$, of entropy $1$, while the whole is
pure, of entropy $0$. Correlations of measurements on entangled pairs violate Bell's inequalities, which any local hidden variable theory satisfies; this was confirmed experimentally (Aspect, Clauser, Zeilinger).
\item \emph{Teleportation and dense coding.} Sharing one Bell pair, two classical bits suffice to transmit one unknown qubit (teleportation), and one qubit suffices to transmit two classical bits (dense coding). The
\emph{Holevo bound} says that $n$ qubits carry at most $n$ classical bits of accessible information.
\item \emph{Quantum error correction.} Although states cannot be copied, they can be protected: Shor's code encodes one qubit in nine and corrects an arbitrary error on any one of them, and the \emph{stabiliser codes} are
built from classical codes over $\mathbb{F}_4$ \cite{nielsen-chuang}. These codes, together with Shor's factoring algorithm, are the theoretical basis of quantum computing (Chapter~\ref{computability-and-complexity}).
\end{enumerate}
\end{remark}

\section{Exercises}\label{information-theory:section-exercises}

\begin{exercise}\label{information-theory:exercise-twelve-coins}
Among $12$ coins one is counterfeit, lighter or heavier. Using entropy, show that at least $3$ weighings on a balance (each with $3$ outcomes) are needed to find it and tell whether it is lighter or heavier, and that $13$ coins
with this requirement would need more than $3$ weighings unless the extra information ``the odd coin is not the thirteenth'' is available.
\end{exercise}

\begin{solution}
There are $24$ equally likely possibilities, so the unknown has entropy $\log 24 \approx 4.585$ bits. Each weighing has $3$ outcomes, so its outcome has entropy at most $\log 3$, and $k$ weighings give at most $k\log 3$ bits of
information (Proposition~\ref{information-theory:proposition-entropy-rules}(3) and (5): the answer must be a function of the outcomes). Since $2\log 3 = \log 9 < \log 24$, two weighings do not suffice. With $13$ coins there are
$26 \leq 27 = 3^3$ possibilities, so counting allows $3$ weighings; but the first weighing of $k$ against $k$ coins leaves $2(13 - 2k)$ possibilities if it balances and $2k$ otherwise, and $\max(26 - 4k, 2k) \leq 9$ is impossible,
so $3$ weighings are not enough without an extra genuine coin.
\end{solution}

\begin{exercise}\label{information-theory:exercise-entropy-function}
On a countable set there is no uniform law and entropy is unbounded (Example~\ref{information-theory:example-entropy}(4)), but a mean constraint bounds it. Prove that among laws on $\{1, 2, \ldots\}$ with $E[X] = \mu > 1$, the geometric law
$P(X = k) = (1 - q)q^{k-1}$ with $q = 1 - \frac1\mu$ has the largest entropy, equal to $\mu h(\frac1\mu)$. In particular, a positive integer valued variable with finite mean has finite entropy.
\end{exercise}

\begin{solution}
Let $g$ be the geometric law. For any $p$ with mean $\mu$, $-\sum_kp(k)\log g(k) = -\log(1 - q) - (\mu - 1)\log q$, which equals $H(g)$ by the same computation for $g$. By Gibbs' inequality, $H(p) \leq -\sum_kp(k)\log g(k) = H(g)$, with
equality iff $p = g$. With $q = 1 - \frac1\mu$, $H(g) = -\log\frac1\mu - (\mu - 1)\log(1 - \frac1\mu) = \mu h(\frac1\mu)$.
\end{solution}

\begin{exercise}\label{information-theory:exercise-huffman-dyadic}
Show that the Huffman code achieves $L = H(X)$ exactly when all probabilities are powers of $\frac12$, and compute its expected length for a fair die.
\end{exercise}

\begin{solution}
By Theorem~\ref{information-theory:theorem-source-coding}(1), $L = H$ requires $p(x) = 2^{-\ell(x)}$. Conversely, for dyadic probabilities the Shannon lengths $\log\frac1{p(x)}$ give $L = H$, and Huffman is optimal, so $L_{\mathrm{Huffman}} \leq H$, hence
$= H$. For a fair die, Huffman merges into lengths $2, 2, 3, 3, 3, 3$ ($\frac16 + \frac16$ three times, then $\frac13 + \frac13$, then $\frac13 + \frac23$), so $L = \frac{2 + 2 + 3 + 3 + 3 + 3}{6} = \frac83 \approx 2.667$ bits, against $H = \log 6 \approx 2.585$.
\end{solution}

\begin{exercise}\label{information-theory:exercise-bsc-repetition}
Over a binary symmetric channel with $p = 0.1$, compare the $3$-fold repetition code (rate $\frac13$, majority decoding) and the $[7, 4]$ Hamming code (rate $\frac47$) by their probabilities of decoding a block wrongly, and compare
both rates with the capacity.
\end{exercise}

\begin{solution}
Repetition fails when at least $2$ of $3$ bits flip: $3p^2(1 - p) + p^3 = 0.028$, per data bit. Hamming decoding fails when at least $2$ of $7$ bits flip: $1 - (1 - p)^7 - 7p(1 - p)^6 \approx 1 - 0.478 - 0.372 = 0.150$ per block of $4$ data
bits. The capacity is $1 - h(0.1) \approx 0.531$, larger than both rates; by Theorem~\ref{information-theory:theorem-channel-coding}, codes of rate $0.5$ with error probability as small as desired exist, but they must be long.
\end{solution}

\begin{exercise}\label{information-theory:exercise-mutual-information-gaussian}
Let $(X, Y)$ be a pair of standard normal variables with correlation $\rho$, $|\rho| < 1$, that is $Y = \rho X + \sqrt{1 - \rho^2}Z$ with $Z \sim N(0, 1)$ independent of $X$. Show that $I(X; Y) = -\frac12\ln(1 - \rho^2)$ nats, defining
$I(X; Y) = h(Y) - h(Y \mid X)$ with $h(Y \mid X) = h(\sqrt{1 - \rho^2}Z)$.
\end{exercise}

\begin{solution}
Given $X$, $Y$ is $N(\rho X, 1 - \rho^2)$, so $h(Y \mid X) = \frac12\ln(2\pi e(1 - \rho^2))$, while $Y \sim N(0, 1)$ has $h(Y) = \frac12\ln(2\pi e)$ (Example~\ref{information-theory:example-differential-entropy}). The difference is
$-\frac12\ln(1 - \rho^2)$, which is $0$ iff $\rho = 0$ and tends to $\infty$ as $|\rho| \to 1$.
\end{solution}
